% TOP_PROBLEM: {"id": 2306101, "title": "Research Problems in Function Theory — Problem 6.101", "category_id": 8, "category": {"id": 8, "name": "analysis", "display_name": "Analysis", "description": "Limits, continuity, calculus, and function theory.", "slug": "analysis", "order_index": 8, "created_at": "2026-07-31T15:26:25.671Z"}, "status": "open"}
% FIRST_POSTED: null
% ENTRY_KIND: statement_only
% Imported from: batch04.tex; UnsolvedMath ID 2306101
\documentclass[11pt]{article}
\usepackage[margin=25mm]{geometry}
\usepackage{fontspec}
\setmainfont{Latin Modern Roman}
\usepackage{amsmath,amssymb,mathtools}
\usepackage{unicode-math}
\setmathfont{Latin Modern Math}
\usepackage{xurl,hyperref}
\hypersetup{colorlinks=true,urlcolor=blue}
\setcounter{secnumdepth}{0}
\setlength{\parindent}{0pt}
\setlength{\parskip}{5pt}
\setlength{\emergencystretch}{3em}
\newcommand{\sref}[2]{\hyperlink{source:#1:#2}{[S#2]}}
\newcommand{\cataloguescope}[1]{\paragraph{Recorded target.}#1}
\begin{document}
% UnsolvedMath ID 2306101; source catalogue code AMR-022-6101
\setcounter{section}{2306100}
\section{Research Problems in Function Theory — Problem 6.101}\label{2306101}
{\small\sffamily UnsolvedMath ID 2306101 · Analysis}
\subsection{Definitions and mathematical statement}

\paragraph{Shared notation.}$\mathbb N=\{1,2,\ldots\}$, and $\mathbb Z,\mathbb Q,\mathbb R,\mathbb C$ denote the integers, rationals, reals and complex numbers. Any other conventions are specified locally.


For a nonconstant holomorphic function $f$ on a domain $D$, its valence is $v_D(f)=\sup_{w\in\mathbb C}n_D(w,f)\in\mathbb N\cup\{\infty\}$, where $n_D(w,f)$ counts the zeros of $f-w$ in $D$ with multiplicity. The function is $p$-valent if $v_D(f)\leq p$. Let $K\subseteq\mathbb C$ be closed and $p\geq1$ an integer. Set
\[\mathcal F(K)=\left\{z\longmapsto\sum_{k=1}^{m}\frac{A_k}{z-a_k}:m\geq1,\ A_k>0,\ a_k\in K\right\}.\]
A common domain of $p$-valence is a nonempty connected open set $D\subseteq\mathbb C\setminus K$ on which $v_D(f)\leq p$ for every $f\in\mathcal F(K)$. Such a domain is maximal if it is contained in no strictly larger common domain of $p$-valence.
\paragraph{Statement recorded in the supplied source.}
For a given closed set $K$ and integer $p>1$, determine a maximal common domain of $p$-valence for $\mathcal F(K)$. The source records that the case $p=1$ was solved by Distler. \sref{2306101}{2}
\subsection{Short English statement}
Find a largest-by-inclusion common region of bounded valence for rational functions with positive residues at poles in a prescribed closed set.
\subsection{Sources}
\begin{description}
\item[\hypertarget{source:2306101:1}{S1}] ULAM.AI and the UnsolvedMath Contributors, UnsolvedMath: A Curated Collection of Open Mathematics Problems, record 2306101 (AMR-022-6101). CC BY 4.0. \url{https://huggingface.co/datasets/ulamai/UnsolvedMath}.
\item[\hypertarget{source:2306101:2}{S2}] W. K. Hayman and E. F. Lingham, Research Problems in Function Theory (2018), Problem 6.101, complete problem and update. \url{https://arxiv.org/abs/1809.07200}.
\end{description}
\label{end:2306101}
\end{document}
